\documentclass[11pt,reqno]{amsart}

\usepackage[T1]{fontenc}
\usepackage[utf8]{inputenc}
\usepackage{lmodern}
\usepackage{amsmath,amssymb,amsthm,mathtools}
\usepackage{enumitem}
\usepackage[margin=1.12in]{geometry}
\usepackage{microtype}
\usepackage[colorlinks=true,linkcolor=blue,citecolor=blue,urlcolor=blue]{hyperref}
\hypersetup{
  pdftitle={A Projective Approach to Non-Dividing Sets},
  pdfsubject={Non-dividing sets, subset sums, and projective density increment},
  pdfkeywords={non-dividing sets, subset sums, generalized arithmetic progressions, convex position, density increment, formal verification}
}

\newtheorem{theorem}{Theorem}[section]
\newtheorem{proposition}[theorem]{Proposition}
\newtheorem{lemma}[theorem]{Lemma}
\newtheorem{corollary}[theorem]{Corollary}
\newtheorem{claim}[theorem]{Claim}
\theoremstyle{definition}
\newtheorem{definition}[theorem]{Definition}
\newtheorem{problem}[theorem]{Problem}
\theoremstyle{remark}
\newtheorem{remark}[theorem]{Remark}

\newcommand{\Z}{\mathbb Z}
\newcommand{\R}{\mathbb R}
\newcommand{\T}{\mathbb T}
\newcommand{\eps}{\varepsilon}
\newcommand{\cA}{\mathcal A}
\newcommand{\cB}{\mathcal B}
\newcommand{\cC}{\mathcal C}
\newcommand{\cH}{\mathcal H}
\newcommand{\cL}{\mathcal L}
\newcommand{\cP}{\mathcal P}
\newcommand{\cS}{\mathcal S}
\newcommand{\cZ}{\mathcal Z}
\newcommand{\Span}{\operatorname{span}}
\newcommand{\conv}{\operatorname{conv}}
\newcommand{\aff}{\operatorname{aff}}
\newcommand{\vol}{\operatorname{vol}}
\newcommand{\diam}{\operatorname{diam}}
\newcommand{\dist}{\operatorname{dist}}
\newcommand{\ind}{\operatorname{ind}}
\newcommand{\rank}{\operatorname{rank}}
\newcommand{\ord}{\operatorname{ord}}
\newcommand{\supp}{\operatorname{supp}}
\newcommand{\SigmaSet}{\Sigma}
\newcommand{\ND}{\mathrm{ND}}
\newcommand{\NA}{\mathrm{NA}}

\title[A projective approach to non-dividing sets]
{A Projective Approach to Non-Dividing Sets}

% Replace ``Anonymous'' by the final author information for a non-anonymous submission.
\author{Theofil Xeff}
\date{}

\begin{document}

\begin{abstract}
Let $F(N)$ be the largest size of a set $A\subseteq\{1,\dots,N\}$ in which
no element divides the sum of any nonempty subset of the remaining
elements.  Erd\H{o}s, Lev, Rauzy, S\'andor and S\'ark\H{o}zy proved
$F(N)=O(N^{1/2})$, and constructions of Erd\H{o}s give
$F(N)\gg N^{1/5}$.  We prove the matching upper bound:
\[
 F(N)=N^{1/5+o(1)}.
\]
The key new idea is a projective normalization $v\mapsto v/\lambda(v)$.
Divisibility constraints depend only on the values of the linear
functional $\lambda$, so they survive the normalization, while the
associated convex geometry drops by one dimension; the density-increment
method behind Pham and Zakharov's sharp non-averaging bound then applies
one dimension lower, moving the exponent from $1/4$ to $1/5$.  The proof
has been formally verified in Lean~4. The novel idea was found by GPT 5.6 Sol.
\end{abstract}

\maketitle

\section{Introduction}
\label{sec:introduction}

For a finite set $A\subseteq\Z_{>0}$, write
\[
 \SigmaSet(A)=\left\{\sum_{a\in S}a:S\subseteq A\right\}
\]
for its set of subset sums.

\begin{definition}
A finite set $A\subseteq\Z_{>0}$ is \emph{non-dividing} if, for every
$a\in A$ and every nonempty set $S\subseteq A\setminus\{a\}$,
\[
 a\nmid \sum_{s\in S}s.
\]
Let $F(N)$ be the maximum cardinality of a non-dividing subset of
$[N]=\{1,\dots,N\}$.
\end{definition}

Erd\H{o}s, Lev, Rauzy, S\'andor and S\'ark\H{o}zy introduced this problem
and proved the bound $F(N)<3N^{1/2}+1$~\cite{ELRSS99}; it appears as
Problem~131 in the collection of Erd\H{o}s problems maintained by
Bloom~\cite{ErdosProblems}.  The previously
known polynomial estimates were
\[
 N^{1/5}\ll F(N)\le N^{1/4+o(1)}.
\]
The lower bound is implicit in the constructions of Erd\H{o}s and of
Erd\H{o}s--Lev--Rauzy--S\'andor--S\'ark\H{o}zy
\cite{Er97,ELRSS99}.  The upper bound follows because every non-dividing
set is non-averaging, together with the sharp theorem of Pham and
Zakharov~\cite{PhamZakharov}.

Our main result closes the gap between these exponents.

\begin{theorem}\label{thm:main}
One has
\[
 F(N)=N^{1/5+o(1)}.
\]
\end{theorem}

\subsection{Outline of the proof}

Every non-dividing set is non-averaging: an averaging relation
$\sum_{s\in S}s=|S|a$ is already forbidden modulo $a$.  The previous
upper bound $F(N)\le N^{1/4+o(1)}$ came from this observation alone,
combined with the sharp non-averaging theorem of
Pham--Zakharov~\cite{PhamZakharov}.  Since the exponent $1/4$ is best
possible for non-averaging sets, any further progress must use the
forbidden divisibilities in an essential way, and the difficulty is
that divisibility, unlike averaging, is destroyed by translation.

Our starting point is that divisibility survives a different family of
operations.  A dyadic pigeonhole places a large part of the set in a
shell $X\le a\le2X$, at the cost of a logarithmic factor.  We then work
in $\Z^d$, with a set $V$ in a symmetric box $B$ and an integer linear
functional $\lambda$ taking values in $[X,2X]$ on $V$, and forbid
$\lambda(a)\mid\sum_{v\in S}\lambda(v)$ for $a\in V$ and nonempty
$S\subseteq V\setminus\{a\}$.  Because these relations see only the
values of $\lambda$, they are preserved exactly when we pass to a subset
or apply the coefficient identification of a proper homogeneous
generalized arithmetic progression (GAP) containing $V$; these two
operations turn out to be all we need.

The structural input is the homogeneous subset-sum theorem of
Conlon--Fox--Pham~\cite{CFPHomogeneous}, in the form recorded in
\cite{PhamZakharov}, which places $V$ in a proper GAP whose subset sums
fill a large translated copy of that GAP.  Iterating coefficient
identifications (Section~\ref{sec:technical-package}), we may assume the
pair $(V,B)$ is \emph{irreducible}: every large subset of $V$ has full
structural dimension, and its structural GAP occupies a fixed proportion
of the box.

The new idea enters here.  Normalizing $\pi(v)=v/\lambda(v)$ sends the
shell into the affine hyperplane $\{\lambda=1\}$, of dimension $d-1$,
and we prove (Theorem~\ref{thm:projective-convexity}) that for an
irreducible pair the projected family $\pi(V)$ must be in approximate
convex position.  Suppose it is not.  A duality argument produces some
$a\in V$ and a bounded positive representation
$Ta=\sum_{v\ne a}\theta_vv$ by the remaining elements; zonotope rounding
converts this identity into genuine subset sums, which fill a
\emph{tube}: every point of a small-index structural lattice in a box
around $Ta$.  The index is small enough that some positive multiple
$m_0a$ lands in the tube, making $m_0\lambda(a)$ a subset sum of the
values $\lambda(v)$ with $v\ne a$ (exactly a forbidden relation).

The rest is a density increment in the style of Pham--Zakharov, run one
dimension lower.  Their convex-density lemma, applied in the hyperplane,
finds a small cap containing many of the points $\pi(v)$, and lifting
the cap back to the lattice (Lemma~\ref{lem:cap-lifting}) shrinks the
box by the volume of the cap.  This gain of one dimension is what moves
the exponent: setting $\sigma_1=\frac15$, $\sigma_2=\frac14$ and
$\sigma_d=\frac{d-2}d$ for $d\ge3$, so that for $d\ge2$ the number
$\sigma_d$ is the sharp non-averaging exponent in dimension $d-1$, the
iteration (Theorem~\ref{thm:projective-main}) shows that a shell
non-dividing set satisfies $|V|\le|B|^{\sigma_d+o(1)}$.  The case $d=1$
is Theorem~\ref{thm:main}.

\subsection{Formal verification}
\label{sec:formalization}

The complete proof of Theorem~\ref{thm:main} has been formalized and
machine-checked in Lean~4 over the mathlib
library~\cite{Formalization}.  The formalization follows the Comparator
workflow of the Lean Focused Research Organization: a short trusted file
states the two definitions and the main theorem, and the kernel certifies
that the development proves exactly that statement.  The only
non-foundational axioms are the statements of the quoted literature
inputs --- the structure theorem of Conlon--Fox--Pham in the form of
\cite[Corollary~5]{PhamZakharov}, the discrete John lemma
\cite[Lemma~7]{PhamZakharov}, zonotope rounding
\cite[Lemma~13]{PhamZakharov}, the convex-density lemma
\cite[Lemma~1]{PhamZakharov}, and the Blaschke selection and volume
continuity theorems from \cite{Schneider} --- so the formalization
verifies precisely the reduction carried out in this paper.  In view of
this, we abbreviate some routine uniformity and bookkeeping verifications
in the longer proofs below; every abbreviated step is carried out in
complete detail in the formalization.

\section{Elementary structure and the lower bound}
\label{sec:elementary}

The bound $F(N)\gg N^{1/5}$ is essentially known, being implicit in the
constructions of \cite{Er97,ELRSS99}; we include a short self-contained
proof.

\subsection{A translation lemma}

A finite set $B\subseteq\Z$ is \emph{non-averaging} if there are no
$b\in B$ and nonempty $S\subseteq B\setminus\{b\}$ such that
\[
 \sum_{s\in S}s=|S|b.
\]

\begin{lemma}
\label{lem:translation}
Let $B\subseteq\Z$ be a non-averaging set of cardinality $m$ and diameter $D$.  If
\[
   x+\min B>mD,
\]
then $x+B$ is non-dividing.
\end{lemma}

\begin{proof}
Suppose that, for $b_j\notin S\subseteq B$,
\[
  x+b_j\mid\sum_{b_i\in S}(x+b_i).
\]
Put $s=|S|$.  Subtracting $s(x+b_j)$, we see that $x+b_j$ divides
\[
  R=\sum_{b_i\in S}(b_i-b_j).
\]
Since $|R|\le sD<mD<x+\min B\le x+b_j$, one has $R=0$.  This is an averaging relation in $B$, a contradiction.
\end{proof}

\subsection{An explicit \texorpdfstring{$N^{1/5}$}{N to the one-fifth} construction}

\begin{proposition}
\label{prop:lower-bound}
For every integer $m\ge2$, there is a non-dividing set of cardinality $m$ contained in $[2m^5]$.  Consequently,
\[
   F(N)\gg N^{1/5}.
\]
\end{proposition}

\begin{proof}
Put
\[
 C=m^3+1,
 \qquad
 b_i=Ci+i^2\quad(1\le i\le m),
\]
and let $B=\{b_1,\dots,b_m\}$.  Suppose that
\[
 \sum_{i\in S}b_i=|S|b_j,
 \qquad j\notin S.
\]
Writing $s=|S|$ gives
\[
 C\left(\sum_{i\in S}i-sj\right)
 +\left(\sum_{i\in S}i^2-sj^2\right)=0.
\]
The second parenthesis is an integer of absolute value at most $m^3<C$.  Reducing the displayed identity modulo $C$ therefore shows that this parenthesis is zero; the first parenthesis is then zero as well.  Hence
\[
 \sum_{i\in S}(i-j)^2=0,
\]
which is impossible because $j\notin S$.  Thus $B$ is non-averaging.

Let $\Delta=b_m-b_1$ and set
\[
  x=m\Delta-b_1+1,
  \qquad
  A=x+B.
\]
Then $\min A=m\Delta+1$, so Lemma~\ref{lem:translation} implies that $A$ is non-dividing.  Moreover,
\[
 \Delta=(m^3+1)(m-1)+m^2-1
       =m^4-m^3+m^2+m-2,
\]
and
\[
 \max A=(m+1)\Delta+1<2m^5
\]
for $m\ge2$.
\end{proof}

\section{The projective framework}
\label{sec:projective-framework}

The upper-bound proof is most naturally formulated in boxes in $\Z^d$.

\begin{definition}[Shell non-dividing set]
Let $B\subseteq\Z^d$ be a symmetric axis-parallel box, let $\lambda:\Z^d\to\Z$ be a nonzero integer linear functional, and let $X>0$.  A finite set $V\subseteq B$ is \emph{$(\lambda,X)$-shell non-dividing} if
\[
 X\le \lambda(v)\le2X
 \qquad(v\in V)
\]
and, for every $a\in V$ and every nonempty $S\subseteq V\setminus\{a\}$,
\[
 \lambda(a)\nmid \sum_{v\in S}\lambda(v).
\]
\end{definition}

The shell restriction loses only a logarithmic factor in the original one-dimensional problem.

\begin{definition}[Projective normalization]
For $v$ in the positive $\lambda$-shell, define
\[
   \pi(v)=\frac{v}{\lambda(v)}.
\]
Then $\pi(v)$ belongs to the affine hyperplane
\[
   H_\lambda=\{y\in\R^d:\lambda(y)=1\},
\]
which has dimension $d-1$.
\end{definition}

This is the source of the expected dimension drop.

\subsection{The exponent sequence}

Define
\[
 \sigma_1=\frac15,
 \qquad
 \sigma_2=\frac14,
 \qquad
 \sigma_d=\frac{d-2}{d}\quad(d\ge3).
\]
For $d\ge2$, this equals the sharp Pham--Zakharov non-averaging exponent in dimension $d-1$.

\section{Structural input and unshifted preprocessing}
\label{sec:technical-package}

We use the higher-dimensional homogeneous subset-sum structure theorem of
Conlon--Fox--Pham~\cite{CFPHomogeneous} in the form recorded by
Pham--Zakharov~\cite[Corollary~5]{PhamZakharov}.  We state only the
consequences required below.

Throughout the technical argument, the ambient dimension and every displayed
polynomial-box exponent are fixed.  Parameters such as $\delta$, $\gamma$ and
$\mu$ may depend on the set size $n$, and all $o(\cdot)$ relations and limits
are taken as $n\to\infty$ along the family under consideration.

We use the following GAP conventions throughout.  A symmetric homogeneous GAP
is data
\[
 P=\left\{\sum_{i=1}^e n_iq_i:|n_i|\le N_i\right\},
 \qquad q_i\in\Z^\ell,\quad N_i\in\Z_{\ge0}.
\]
Its \emph{dimension} is the number $e$ of displayed steps, while its
\emph{rank} is $\dim_{\R}\Span_{\R}\{q_1,\dots,q_e\}$.  It is
\emph{proper} if its coefficient map is injective on the displayed box.  For
$t\ge0$ we set
\[
 tP=\left\{\sum_{i=1}^e n_iq_i:|n_i|\le\lfloor tN_i\rfloor\right\}.
\]
If $P$ is proper, $\phi_P$ denotes the inverse coefficient identification from
$P$ to its coordinate box.  This is an identification on $P$, not in general
an ambient homomorphism $\Z^\ell\to\Z^e$.  For a real number $x\ge3$ we also
write
\[
 s(x)=\frac{x}{(\log x)^2};
\]
thus $s(C)=s(|C|)$ below.

\begin{theorem}[Conlon--Fox--Pham; Pham--Zakharov]
\label{thm:cfp-structure}
For every fixed ambient dimension $\ell$ and every $\beta>1$ there are constants $c>0$ and $D\in\mathbb N$ with the following property.  Let $C\subseteq\Z^\ell$ have cardinality $m$ and be contained in an axis-parallel box $Q$ satisfying $|Q|\le m^\beta$.  For all sufficiently large $m$ there are a set $\widehat C\subseteq C$ with
\[
 |\widehat C|\ge m-c^{-1}\frac{m}{\log m},
\]
a symmetric GAP
\[
 P(C)=\left\{\sum_{i=1}^{d'}n_iq_i:|n_i|\le N_i\right\},
 \qquad d'\le D,
\]
containing $\widehat C\cup\{0\}$, and a set $C_0\subseteq\widehat C$ with $|C_0|\le s(C)$ such that
\[
 q+c s(C)P(C)\subseteq\SigmaSet(C_0)
\]
for some $q\in s(C)P(C)$.  Moreover $c s(C)P(C)$ is proper.
\end{theorem}

The existence statement through properness is the cited corollary of Pham--Zakharov, derived from the Conlon--Fox--Pham structure theorem.  Their symmetric convention is the centered display above, so $P(C)$ is homogeneous.  The location $q\in s(C)P(C)$ is not an extra external hypothesis: since $C_0\subseteq\widehat C\subseteq P(C)$ and $|C_0|\le s(C)$, one has $\SigmaSet(C_0)\subseteq s(C)P(C)$; since $0\in c s(C)P(C)$, the core inclusion itself gives $q\in\SigmaSet(C_0)\subseteq s(C)P(C)$.

For fixed ambient dimension and polynomial-box exponent, a
\emph{structural package} means a fixed choice of the constants and of the
objects $\widehat C,C_0,q,P(C)$ in Theorem~\ref{thm:cfp-structure} for every
eligible sufficiently large set $C$.  When dimensions vary over a fixed
finite range, we choose one package in each ambient dimension before
beginning an iteration and, after every coefficient identification, use
the package indexed by the new ambient dimension; no post hoc choice of a
polynomial-box exponent or of structural constants is ever made, so all
constants are uniform once the finite dimension range and the common
polynomial-box exponent have been fixed.  Within each package we choose
an admissible structural GAP of minimum dimension and write
$d(C)=\dim P(C)$; ties are broken arbitrarily, and every radius then
satisfies $N_i\ge1$.  For sufficiently large $|C|$, properness of
$cs(C)P(C)$ implies properness of $P(C)$, and the centered representation
makes $\phi_{P(C)}$ based at the origin.  We suppress the chosen package
from the notation.

\begin{definition}[Unshifted irreducibility]
Fix a structural package applicable to all sets under consideration.  Let
$V\subseteq B\subseteq\Z^d$, where $B$ is a symmetric axis-parallel
coordinate box.  We say that the pair $(V,B)$ is
\emph{$(\delta,\gamma)$-irreducible in the unshifted sense}, relative to
this package, if every $D\subseteq V$ with $|D|\ge\delta|V|$ has
subset-sum dimension $d$ and satisfies
\[
 |P(D)|\ge\gamma|B|.
\]
\end{definition}

The omission of translations is essential: shell non-divisibility is inherited
by subsets and by coefficient identifications of proper homogeneous GAPs, with
the shell functional transported as below, but not by arbitrary translations.

\begin{lemma}[Transformation of the shell functional]
\label{lem:functional-transform}
Let $P\subseteq\Z^\ell$ be a proper homogeneous GAP with step vectors $q_1,\dots,q_d$, and let $\lambda:\Z^\ell\to\Z$ be integer-linear.  Define
\[
 \lambda_P(n_1,\dots,n_d)=\sum_{i=1}^d n_i\lambda(q_i).
\]
Then
\[
 \lambda_P(\phi_P(v))=\lambda(v)\qquad(v\in P).
\]
Consequently, if $V\subseteq P$ is shell non-dividing for $\lambda$, then $\phi_P(V)$ is shell non-dividing for $\lambda_P$ with exactly the same shell parameter.
\end{lemma}

\begin{proof}
Homogeneity gives $v=\sum_i n_iq_i$ whenever $\phi_P(v)=(n_1,\dots,n_d)$.  The identity follows by linearity.  Every prohibited divisibility relation is an identity involving only the integer values of the functional, so it is preserved exactly.
\end{proof}

\begin{lemma}[Progression comparison in a coordinate box]
\label{lem:progression-comparison}
Fix $r$ and $\beta>1$.  Let $Q\subseteq\Z^r$ be a proper symmetric
coordinate GAP and let $C\subseteq Q$ satisfy $|Q|\le|C|^\beta$.  Write
$e=d(C)$ and $s=s(C)$, using the structural package for $(r,\beta)$.  Then
\[
 |P(C)|\ll_{r,\beta} s^{-(e-r)}|Q|\qquad(e>r),
\]
while
\[
 |P(C)|\ll_{r,\beta} |Q|\qquad(e\le r).
\]
The constants are uniform when $r$, $\beta$ and the structural dimensions
under consideration range over fixed bounded sets.
\end{lemma}

\begin{proof}
Let $R=P(C)$, and let $C_0\subseteq C$ be the core set supplied by Theorem~\ref{thm:cfp-structure}.  Since $C_0\subseteq Q$ and $|C_0|\le s$,
\[
 q+csR\subseteq\SigmaSet(C_0)\subseteq sQ.
\]
If $e>r$, properness of $csR$ and of the coordinate GAP $Q$ gives
\[
 s^e|R|\ll |csR|\le |sQ|\ll_r s^r|Q|,
\]
which proves the first assertion.

Suppose $e\le r$.  Let $H$ be the real span of the steps of $R$, of
dimension $h\le e$; for all sufficiently large $|C|$ the affine hull of
$q+csR$ has direction space $H$.  Differences of points of
$q+csR\subseteq\SigmaSet(C_0)\subseteq\Span_{\R}(C_0)$ show that
$H\subseteq\Span_{\R}(C_0)$; conversely $C_0\subseteq R\subseteq H$, so
$C_0$ spans $H$ and, in particular, the lattice points of
$H\cap\conv(Q)$ span $H$.  Apply the fixed-dimensional discrete John
lemma \cite[Lemma~7]{PhamZakharov} to the symmetric convex body
$H\cap\conv(Q)$ in the lattice $H\cap\Z^r$, identified with $\Z^h$.
Since $Q$ is a coordinate box,
$(H\cap\conv(Q))\cap\Z^r=H\cap Q$, so the lemma gives a proper symmetric
rank-$h$ GAP $Q_H$ with
\[
 H\cap Q\subseteq Q_H,
 \qquad
 |Q_H|\ll_r |H\cap Q|\le |Q|.
\]
The set $C_0$ lies in $R\cap Q\subseteq H\cap Q$, hence
$q+csR\subseteq\SigmaSet(C_0)\subseteq sQ_H$, and consequently
\[
 s^e|R|\ll |csR|
 \le |sQ_H|
 \ll_r s^h|Q_H|
 \ll_r s^e|Q|,
\]
proving the second assertion.
\end{proof}

\begin{lemma}[Unshifted preprocessing]
\label{lem:unshifted-preprocessing}
Fix an ambient dimension $\ell$, $\beta>1$, $\eps\in(0,1/3)$ and a
desired shrink exponent $K_0\ge1$.  Let
$V\subseteq B\subseteq\Z^\ell$ be shell non-dividing, with
$|B|\le |V|^\beta$.  Suppose $0<\delta\le1/2$ and
\[
 \gamma\le\delta^K,
 \qquad
 \gamma\ge |V|^{-1/3},
\]
where $K=K(\ell,\beta,\eps,K_0)$ is sufficiently large.  Then one can
pass, using only subsets and coefficient identifications of proper
homogeneous GAPs, to a shell non-dividing set
$\widetilde V\subseteq\widetilde B\subseteq\Z^{\widetilde d}$ and
fix a structural package in $\Z^{\widetilde d}$ such that
$(\widetilde V,\widetilde B)$ is $(\delta,\gamma)$-irreducible relative
to that package and
\begin{align*}
 |\widetilde B|&\ll_{\ell,\beta,\eps,K_0} |V|^{-(1-\eps)(\widetilde d-\ell)}|B|,
 &|\widetilde V|&>|V|^{1-\eps},&&\widetilde d>\ell,\\
 |\widetilde B|&\ll_{\ell,\beta,\eps,K_0} \rho^{K_0}|B|,
 &|\widetilde V|&=\rho|V|>|V|^{1-\eps},&&\widetilde d=\ell,\\
 |\widetilde B|&\ll_{\ell,\beta,\eps,K_0}|B|,
 &|\widetilde V|&>|V|^{1-\eps},&&\widetilde d<\ell.
\end{align*}
Here $\widetilde d=O_{\ell,\beta}(1)$; this dimension bound is uniform
for $\eps\in(0,1/3)$ and $K_0\ge1$.  The polynomial-box exponent and
constants of the output structural package are bounded in
$\ell,\beta,\eps,K_0$.
\end{lemma}

\begin{proof}
Put $n=|V|$ and $\delta_*=\delta/2$; we assume throughout that $n$ is
sufficiently large in the fixed parameters and enlarge $K$ whenever
needed.  Let $D_{\rm init}=D(\ell,\beta)$ be the structural-dimension
bound in Theorem~\ref{thm:cfp-structure}, and choose
\[
 C_U>3(\beta+2),
 \qquad
 \bar d>D_{\rm init}+C_U,
 \qquad
 \bar\beta=4(\beta+3).
\]
Before beginning the iteration, fix the initial package for
$(\ell,\beta)$, structural packages with common polynomial-box exponent
$\bar\beta$ in every ambient dimension $1\le r\le\bar d$, the bound
$D_*=\max_{1\le r\le\bar d}D(r,\bar\beta)$, and a constant
$C=C(\bar d,\bar\beta)$ dominating the constants of
Lemma~\ref{lem:progression-comparison} whenever the ambient dimension is
at most $\bar d$ and the output structural dimension is at most $D_*$.

Apply the structural package for $(\ell,\beta)$ to $V\subseteq B$, and write
$A_0=V$, $P_0=P(V)$ and $d_0=\dim P_0$.  Lemma~\ref{lem:progression-comparison},
applied initially in $B$, and $s(V)\ge n^{1-\eps}$ give
\begin{equation}
\label{eq:initial-progression-bound}
 |P_0|\ll_{\ell,\beta}
 n^{-(1-\eps)\max(0,d_0-\ell)}|B|
 \le n^{\beta+1}
\end{equation}
for all sufficiently large $n$.

At an accepted stage $A_j$ let $P_j$ be its assigned structural progression
of dimension $d_j$, and define
\[
 A_j^\sharp=\phi_{P_j}(\widehat A_j),
 \qquad Q_j=\phi_{P_j}(P_j),
\]
so that $Q_j$ is a proper symmetric coordinate box of dimension $d_j$ and
cardinality $|P_j|$, while
$|A_j^\sharp|=(1-o(1))|A_j|\ge\frac12|A_j|$.  We impose the provisional
eligibility condition
\begin{equation}
\label{eq:package-eligibility}
 |Q_j|\le(\delta|A_j^\sharp|)^{\bar\beta};
\end{equation}
if it fails the procedure stops, and if it holds every
$D\subseteq A_j^\sharp$ with $|D|\ge\delta|A_j^\sharp|$ is eligible for the
already fixed package in ambient dimension $d_j$.

If $(A_j^\sharp,Q_j)$ is not unshiftedly
$(\delta,\gamma)$-irreducible, choose a witnessing set
$A_{j+1}\subseteq A_j^\sharp$ with
$|A_{j+1}|\ge\delta|A_j^\sharp|\ge\delta_*|A_j|$, its structural data
taken from the fixed package in dimension $d_j$.  We call the proposed
move a down-move if $d_{j+1}<d_j$, an up-move if $d_{j+1}>d_j$, and a
shrink-move if $d_{j+1}=d_j$ and $|P_{j+1}|<\gamma|P_j|$.
Lemma~\ref{lem:progression-comparison} gives
\begin{align*}
 |P_{j+1}|&\le C|P_j|&&\text{for a down-move},\\
 |P_{j+1}|&\le C s(A_{j+1})^{-(d_{j+1}-d_j)}|P_j|
     &&\text{for an up-move},\\
 |P_{j+1}|&<\gamma|P_j|&&\text{for a shrink-move}.
\end{align*}
At each coefficient identification the shell functional is transported by
Lemma~\ref{lem:functional-transform}; hence every accepted set remains shell
non-dividing.

Accept proposed moves sequentially, but stop before accepting one if its
cardinality is at most $n^{1-\eps/2}$ or if its structural dimension
exceeds $\bar d$; otherwise stop when the eligibility condition fails or
no move is available.  Suppose that $A_0,\dots,A_i$ have been accepted,
let $\mathcal D,\mathcal U,\mathcal S$ be the sets of down-, up- and
shrink-move indices, and put $U=\sum_{j\in\mathcal U}(d_{j+1}-d_j)$.
Multiplying the move estimates gives
\begin{equation}
\label{eq:unshifted-product}
 |P_i|
 \le C^{|\mathcal D|+|\mathcal U|}
 \gamma^{|\mathcal S|}
 \prod_{j\in\mathcal U}s(A_{j+1})^{-(d_{j+1}-d_j)}|P_0|,
\end{equation}
while the dimension bookkeeping gives $d_i-d_0\le U-|\mathcal D|$, hence
$|\mathcal D|\le d_0+U$ and $|\mathcal U|\le U$; moreover
$|A_i|\ge\delta_*^in$.  Every accepted set after $A_0$ has size greater
than $n^{1-\eps/2}$, so $s(A_j)\ge n^{1-\eps}$ for every $j\ge1$
occurring in an up-move.  Since $|P_i|\ge1$ and $\gamma\le1$,
equation~\eqref{eq:unshifted-product} together
with~\eqref{eq:initial-progression-bound} implies, once
$C^{D_{\rm init}}\le n$ and $C^2\le n^{(1-\eps)/2}$,
\[
 n^{(1-\eps)U/2}\le n^{\beta+2},
 \qquad\text{and hence}\qquad
 U\le\frac{2(\beta+2)}{1-\eps}
 \le3(\beta+2)<C_U.
\]

We check that none of the three provisional cutoffs ever triggers; the
uniformity bookkeeping abbreviated here is carried out in full in the
formalization (Section~\ref{sec:formalization}).

\emph{Dimension.}  Every accepted dimension is at most
$d_0+U<D_{\rm init}+C_U<\bar d$.  The identical calculation applies after
adjoining any proposed move whose cardinality exceeds $n^{1-\eps/2}$: its
structural dimension is at most $D_*$, so the same constant $C$ applies,
the total upward change is still less than $C_U$, and the proposed
dimension is also less than $\bar d$.

\emph{Eligibility.}  The same estimates applied to every partial
sequence, discarding all factors at most one
in~\eqref{eq:unshifted-product}, give
\begin{equation}
\label{eq:coarse-partial-box}
 |P_i|\le C^{D_{\rm init}+2C_U}|P_0|\le n^{\beta+2}.
\end{equation}
On the other hand, $K\ge1$, $\gamma\ge n^{-1/3}$ and $|A_i|>n^{1-\eps/2}$
give
$\delta|A_i^\sharp|\ge\frac12n^{1-\eps/2-1/(3K)}\ge\frac12n^{2/3-\eps/2}$;
since $\eps<1/3$ and $\bar\beta=4(\beta+3)$, the right-hand side raised
to the power $\bar\beta$ exceeds $n^{\beta+2}$.  Together
with~\eqref{eq:coarse-partial-box}, this
proves~\eqref{eq:package-eligibility} at every accepted stage, so the
entire argument uses only the packages fixed before the iteration.

\emph{Cardinality.}  Equation~\eqref{eq:unshifted-product} and the bound
on $U$ give
\begin{equation}
\label{eq:unshifted-product-simplified}
 |P_i|\ll_{\ell,\beta,\eps}
 \gamma^{i-|\mathcal U|-|\mathcal D|}
 n^{-(1-\eps)U}|P_0|.
\end{equation}
A proposed set $A_{i+1}$ of size at most $n^{1-\eps/2}$ would satisfy
$\delta_*^{i+1}n\le|A_{i+1}|\le n^{1-\eps/2}$, hence
$\delta_*^{i+1}\le n^{-\eps/2}$.  Since $\delta\le1/2$, one has
$\delta^2\le\delta_*$ and therefore
$\gamma\le\delta^K\le\delta_*^{K/2}$.  With
$C_1=C_1(\ell,\beta,\eps)$ bounding $|\mathcal U|+|\mathcal D|$,
$|P_0|\le n^{\beta+1}$ and $\gamma\ge n^{-1/3}$, we get
\[
 \gamma^{i-|\mathcal U|-|\mathcal D|}|P_0|
 \le\gamma^{i+1}\gamma^{-(C_1+1)}|P_0|
 \le n^{-\eps K/4+(C_1+1)/3+\beta+1}.
\]
Taking $K$ sufficiently large in $\ell,\beta,\eps,K_0$ and then $n$
sufficiently large makes this smaller than the reciprocal of the implicit
constant in~\eqref{eq:unshifted-product-simplified}, contradicting
$|P_i|\ge1$.

The procedure therefore stops only because no move is available: every
$D\subseteq A_i^\sharp$ with $|D|\ge\delta|A_i^\sharp|$ has $d(D)=d_i$
and $|P(D)|\ge\gamma|P_i|$, so $(A_i^\sharp,Q_i)$ is unshiftedly
$(\delta,\gamma)$-irreducible relative to the fixed package in dimension
$d_i$.  Put
\[
 \widetilde V=A_i^\sharp,
 \qquad
 \widetilde B=Q_i,
 \qquad
 \widetilde d=d_i.
\]
The deletion estimate and the stopping rule give
$|\widetilde V|>|V|^{1-\eps}$, and the terminal package belongs to the
preselected finite collection, whose dimensions, exponent $\bar\beta$,
constants and thresholds depend only on $\ell,\beta,\eps,K_0$.

It remains to prove the three box estimates.  Combining
\eqref{eq:unshifted-product-simplified} with
\eqref{eq:initial-progression-bound} yields
\begin{equation}
\label{eq:final-box-master}
 |P_i|\ll_{\ell,\beta,\eps}
 \gamma^{i-|\mathcal U|-|\mathcal D|}
 n^{-(1-\eps)(U+\max(0,d_0-\ell))}|B|.
\end{equation}
If $d_i>\ell$, then $U+\max(0,d_0-\ell)\ge d_i-\ell$,
so~\eqref{eq:final-box-master} gives the first displayed alternative; if
$d_i<\ell$, discarding all factors at most one gives
$|P_i|\ll_{\ell,\beta,\eps}|B|$.  Finally suppose $d_i=\ell$ and put
$M=U+\max(0,d_0-\ell)$, so that the dimension bookkeeping gives
$|\mathcal U|+|\mathcal D|\le2M$.  Since $\eps<1/3$ and
$\gamma\ge n^{-1/3}$, one has $n^{-(1-\eps)M}\le\gamma^{2M}$, and
equation~\eqref{eq:final-box-master} gives
$|P_i|\ll_{\ell,\beta,\eps}\gamma^i|B|$.  Furthermore
$|A_i|\ge\delta_*^in$ and $|\widetilde V|=(1-o(1))|A_i|$, so
$\rho=|\widetilde V|/n$ satisfies $\delta_*^i\ll\rho$; since
$\gamma\le\delta_*^{K/2}$, choosing $K\ge2K_0$ sufficiently large gives
$\gamma^i\ll_{K_0}\rho^{K_0}$.  This proves the same-dimensional
alternative and completes the proof.
\end{proof}


\section{Three geometric--arithmetic lemmas}
\label{sec:three-lemmas}

We next establish the slab, rounding and lattice-index statements used in the tube argument.

\subsection{A slab dichotomy}

For a symmetric real box $B_{\R}=\prod_{i=1}^d[-N_i,N_i]$ with
$N_i\in\Z_{\ge1}$, write $B=B_{\R}\cap\Z^d$ and
\[
 w_B(u)=\max_{x\in B_{\R}}|u(x)|.
\]

\begin{lemma}[Convex lattice sets and their John progression]
\label{lem:convex-lattice-john}
Let $H\subseteq\R^d$ be an $r$-dimensional rational subspace, let
$\Lambda=H\cap\Z^d$, and let $K\subseteq H$ be an origin-symmetric convex
body such that $K\cap\Lambda$ spans $H$.  Then $K\cap\Lambda$ is contained
in a proper symmetric rank-$r$ GAP $Q\subseteq\Lambda$ satisfying
\[
 |Q|\ll_d |K\cap\Lambda|.
\]
If $r=d$, then additionally
\[
 |K\cap\Z^d|\ll_d \vol_d(K).
\]
\end{lemma}

\begin{proof}
The case $r=0$ is immediate.  For $r\ge1$, identify $\Lambda$ with
$\Z^r$ and apply the discrete John lemma \cite[Lemma~7]{PhamZakharov}.
It gives a lattice basis $v_1,\dots,v_r$, radii $N_i\ge1$, and a constant
$c_r>0$ such that
\[
 \left\{\sum n_iv_i:|n_i|\le\lfloor c_rN_i\rfloor\right\}
 \subseteq K\cap\Lambda
 \subseteq Q:=\left\{\sum n_iv_i:|n_i|\le N_i\right\}.
\]
The outer progression is proper and $|Q|\ll_r|K\cap\Lambda|$.  If $r=d$,
the successive-minima lattice-point bound and Minkowski's second theorem give
\[
 |K\cap\Z^d|\ll_d\prod_{i=1}^d\lambda_i^{-1}\ll_d\vol_d(K),
\]
since the last successive minimum is at most one.
\end{proof}

\begin{lemma}[Lattice points in a central slab]
\label{lem:slab-count}
Let $u$ be a nonzero real linear functional, let $0<\tau\le1$, and put
\[
 K=B_{\R}\cap\{x:|u(x)|\le\tau w_B(u)\},
 \qquad
 S=K\cap\Z^d.
\]
If the real span of $S$ has dimension $r$, then $S$ is contained in a
proper symmetric rank-$r$ GAP $Q$ satisfying
\[
 |Q|\ll_d
 \begin{cases}
  \tau|B|,&r=d,\\
  |B|,&r<d.
 \end{cases}
\]
\end{lemma}

\begin{proof}
Let $H=\Span_{\R}S$ and $\Lambda=H\cap\Z^d$.  The set $S$ is precisely the
lattice-point set of the origin-symmetric convex body $K\cap H$ in
$\Lambda$, and it spans $H$.  Lemma~\ref{lem:convex-lattice-john} therefore
places $S$ in a proper rank-$r$ GAP $Q$ with
\[
 |Q|\ll_d|S|.
\]
If $r<d$, the trivial estimate $|S|\le|B|$ proves the second alternative.

Assume $r=d$.  Lemma~\ref{lem:convex-lattice-john} gives
\[
 |S|\ll_d\vol_d(K).
\]
Write $u(x)=\sum_i a_ix_i$ and choose an index $j$ for which
$|a_j|N_j\ge w_B(u)/d$.  After fixing all coordinates other than $x_j$,
the allowed values of $x_j$ form an interval whose length is at most
\[
 \frac{2\tau w_B(u)}{|a_j|}\le2d\tau N_j.
\]
Fubini's theorem consequently gives
\[
 \vol_d(K)\le d\tau\vol_d(B_{\R}).
\]
Since the side widths of the coordinate box are positive integers,
$\vol_d(B_{\R})\asymp_d|B|$.  Combining the last three displays proves
$|Q|\ll_d\tau|B|$.
\end{proof}

\begin{lemma}[Slab exclusion from irreducibility]
\label{lem:slab-exclusion}
Fix $d$ and $\beta>1$.  Let $V\subseteq B\subseteq\Z^d$, and fix a
structural package arising from Theorem~\ref{thm:cfp-structure} in ambient
dimension $d$ with polynomial-box exponent at most $\beta$, applicable to
every subset of $V$ of size at least $\delta|V|$.  Suppose that $(V,B)$ is
unshiftedly $(\delta,\gamma)$-irreducible relative to this package, and put
$n=|V|$, and assume
\[
 |B|\le n^\beta,
 \qquad
 \delta\ge n^{-1/2},
 \qquad
 [s(\delta n)]^{-1}=o(\gamma).
\]
There is $c_{d,\beta}>0$ such that every central slab of relative width at
most $c_{d,\beta}\gamma$ contains fewer than $\delta n$ points of $V$.
\end{lemma}

\begin{proof}
Suppose that $D\subseteq V$ lies in a central slab of relative width
$\tau\le c_{d,\beta}\gamma$ and that $|D|\ge\delta n$.  Since
$\delta\ge n^{-1/2}$, one has $|B|\le n^\beta\le|D|^{2\beta}$, so the
fixed structural package is applicable to $D$, and irreducibility gives
$d(D)=d$ and $|P(D)|\ge\gamma|B|$.

Let $r$ be the dimension of the real span of all lattice points in the
slab, let $Q$ be supplied by Lemma~\ref{lem:slab-count}, and put
$s_D=s(D)$.  The structural core $D_0\subseteq D\subseteq Q$ has size at
most $s_D$, so
$q+cs_DP(D)\subseteq\SigmaSet(D_0)\subseteq s_DQ$, and the left-hand
GAP is proper.  Consequently
\[
 s_D^d|P(D)|\ll_{d,\beta}|cs_DP(D)|
 \le |s_DQ|
 \ll_d s_D^r|Q|.
\]
If $r<d$, then
$|P(D)|\ll_{d,\beta}s_D^{r-d}|Q|\le C_{d,\beta}s_D^{-1}|B|$; since
$x/(\log x)^2$ is increasing for large $x$, one has
$s_D\ge s(\delta n)$, and the parameter hypothesis makes the last
expression $o(\gamma|B|)$, a contradiction.  If $r=d$, then
\[
 |P(D)|\ll_{d,\beta}|Q|
 \ll_d\tau|B|
 \le C_{d,\beta}c_{d,\beta}\gamma|B|,
\]
and choosing $c_{d,\beta}$ sufficiently small gives the final
contradiction.
\end{proof}

\subsection{One-sided filling around a zonotope}

For a finite set $E\subseteq\R^d$, write
\[
 Z(E)=\left\{\sum_{v\in E}t_vv:0\le t_v\le1\right\}
\]
for its zonotope.  For a homogeneous GAP $P$, let $\langle P\rangle$ be the lattice generated by its steps.

\begin{lemma}[Rounding with a homogeneous core]
\label{lem:one-sided-filling}
Let $C\subseteq B\subseteq\Z^d$, let $P=P(C)$ have dimension $d$, and assume that its step lattice
\[
 L=\langle P\rangle
\]
has rank $d$.  Take the objects $\widehat C,C_0,q$ from Theorem~\ref{thm:cfp-structure}, and put
\[
 \overline C=\widehat C\setminus C_0,
 \qquad s=s(C).
\]
For all sufficiently large $|C|$,
\[
 L\cap\left(Z(\overline C)+q+\frac{cs}{2}P\right)
 \subseteq\SigmaSet(C).
\]
\end{lemma}

\begin{proof}
Write $P=\bigl\{\sum_{i=1}^d n_iq_i:|n_i|\le N_i\bigr\}$.  The rank
assumption and the fact that $P$ has exactly $d$ steps imply that
$q_1,\dots,q_d$ are linearly independent, hence a $\Z$-basis of $L$, and
the coefficient map extends to a real-linear isomorphism
$\psi:\operatorname{span}_{\R}L\to\R^d$,
$\psi(\sum_i x_iq_i)=(x_1,\dots,x_d)$, with
$\psi(P)=\prod_{i=1}^d[-N_i,N_i]\cap\Z^d$.

Let
\[
 y=z+t\in L\cap\left(Z(\overline C)+q+\frac{cs}{2}P\right),
\]
where $z\in Z(\overline C)$ and $t\in q+(cs/2)P$.  Since $q\in sP$, both
$q$ and $t$ belong to $L$; hence $z=y-t\in L$ and
$\psi(z)\in Z(\psi(\overline C))$.  Apply the Conlon--Fox--Pham
zonotope-rounding lemma in the form of \cite[Lemma~13]{PhamZakharov} to
$A^*=\psi(\overline C)\cup\{0\}$ inside the coefficient box
$\prod_i[-N_i,N_i]\cap\Z^d$, whose $i$th width is $2N_i+1$ in the
convention of that lemma.  It gives a subset sum
$\widetilde z'\in\SigmaSet(\psi(\overline C))$ such that
\[
 |\psi(z)_i-\widetilde z'_i|
 \le \sqrt{d|A^*|}\,(2N_i+1)
 \le 3\sqrt{2d|C|}\,N_i
 \qquad(1\le i\le d),
\]
using $|A^*|\le|C|+1\le2|C|$ and $N_i\ge1$; the rounding error carries
the sign of \cite[equation~(15)]{PhamZakharov}, whose rounding is written
$z=\widetilde z+r$.  Pulling back by $\psi$ gives a subset sum
$\widetilde z\in\SigmaSet(\overline C)$ with
\[
 z-\widetilde z\in C_d\sqrt{|C|}\,P,
 \qquad
 C_d=3\sqrt{2d}.
\]
Since $s=|C|/(\log|C|)^2$, one has $C_d\sqrt{|C|}\le cs/2$ for all
sufficiently large $|C|$.  As $t\in q+(cs/2)P$, coefficientwise addition
in the symmetric GAP therefore gives
\[
 (z-\widetilde z)+t\in q+csP\subseteq\SigmaSet(C_0).
\]
The subset producing $\widetilde z$ lies in $\overline C$ and the subset
producing $(z-\widetilde z)+t$ lies in $C_0$; these are disjoint, so
$y=\widetilde z+\bigl((z-\widetilde z)+t\bigr)\in\SigmaSet(C)$.
\end{proof}

\subsection{A large progression has small lattice index}

\begin{lemma}[Containment of the structural GAP]
\label{lem:gap-contained}
Let $C\subseteq B\subseteq\Z^d$, where $B=B_{\R}\cap\Z^d$ and $B_{\R}$ is a symmetric axis-parallel real box.  Let $P=P(C)$ and let $c$ be the constant in Theorem~\ref{thm:cfp-structure}.  For all sufficiently large $|C|$,
\[
 P\subseteq \frac{4}{c}B_{\R}.
\]
In particular, when the ambient dimension and polynomial-box exponent are fixed, $P\subseteq C_{d,\beta}B_{\R}$.
\end{lemma}

\begin{proof}
Let $s=s(C)$ and let $C_0,q$ be the structural core and translation.  Since
$C_0\subseteq B$ and $|C_0|\le s$,
\[
 q+csP\subseteq\SigmaSet(C_0)\subseteq sB_{\R}.
\]
As $0\in csP$, also $q\in sB_{\R}$, and hence $csP\subseteq2sB_{\R}$.
For $p\in P$, put $k=\lfloor cs\rfloor\ge cs/2$.  Then $kp\in csP$, so
$p\in(2s/k)B_{\R}\subseteq(4/c)B_{\R}$.
\end{proof}

\begin{lemma}[Lattice sections of a coordinate box]
\label{lem:lattice-section-box}
Let
\[
 B_{\R}=\prod_{i=1}^d[-M_i,M_i],
 \qquad
 B=B_{\R}\cap\Z^d,
\]
where the $M_i$ are positive integers.  If $H\subseteq\R^d$ is an
$r$-dimensional linear subspace and $t\ge1$, then
\[
 |H\cap tB_{\R}\cap\Z^d|\le t^r|B|.
\]
\end{lemma}

\begin{proof}
Choose $I\subseteq[d]$, $|I|=r$, so that the coordinate projection
$\operatorname{pr}_I:H\to\R^I$ is injective.  It is then injective on
$H\cap\Z^d$, and therefore
\[
 |H\cap tB_{\R}\cap\Z^d|
 \le\prod_{i\in I}(2\lfloor tM_i\rfloor+1)
 \le t^r\prod_{i\in I}(2M_i+1)\le t^r|B|.
\]
\end{proof}

\begin{lemma}[Full rank and index]
\label{lem:index-bound}
Fix $d$ and $\beta>1$, and fix a structural package arising from
Theorem~\ref{thm:cfp-structure} in ambient dimension $d$ with polynomial-box
exponent at most $\beta$.  Let $(V,B)$ be unshiftedly
$(\delta,\gamma)$-irreducible relative to this package, where
$B=B_{\R}\cap\Z^d$ is a symmetric coordinate box.  Let $C\subseteq V$
have $|C|\ge\delta|V|$ and assume $|B|\le|C|^\beta$.  Put
\[
 P=P(C),\qquad s=s(C),\qquad L=\langle P\rangle.
\]
If
\[
 s\gamma\longrightarrow\infty,
\]
then $L$ has rank $d$ and
\[
 [\Z^d:L]\ll_{d,\beta}\gamma^{-1}.
\]
\end{lemma}

\begin{proof}
Irreducibility gives $\dim P=d$ and $|P|\ge\gamma|B|$.  Let
$c=c(d,\beta)>0$ be the structural constant in
Theorem~\ref{thm:cfp-structure}; the proof of
Lemma~\ref{lem:gap-contained} gives the stronger core containment
\begin{equation}
 csP\subseteq2sB_{\R}.
 \label{eq:index-core-containment}
\end{equation}
Write $P=\bigl\{\sum_{i=1}^d n_iq_i:|n_i|\le N_i\bigr\}$, where every
$N_i\ge1$ as zero-width directions are not counted in the dimension, and
let $H=\Span_{\R}\{q_1,\dots,q_d\}$, of dimension $r$.  Because $csP$ is
proper, its cardinality is the product of its coefficient widths, so for
all sufficiently large $s$
\begin{equation}
 |csP|
 =\prod_{i=1}^d\bigl(2\lfloor csN_i\rfloor+1\bigr)
 \gg_{d,\beta}s^d|P|,
 \label{eq:index-proper-lower}
\end{equation}
while \eqref{eq:index-core-containment} and
Lemma~\ref{lem:lattice-section-box} give
\begin{equation}
 |csP|
 \le |H\cap2sB_{\R}\cap\Z^d|
 \le (2s)^r|B|.
 \label{eq:index-section-upper}
\end{equation}
Combining \eqref{eq:index-proper-lower} and
\eqref{eq:index-section-upper} with $|P|\ge\gamma|B|$ yields
$s^{d-r}\gamma\ll_{d,\beta}1$.  If $r<d$, the left-hand side is at least
$s\gamma$ for large $s$, contradicting the hypothesis.  Hence $r=d$, the
$d$ step vectors are linearly independent, and they form a $\Z$-basis of
their integer span $L$.

Let $M$ be the $d\times d$ matrix whose columns are $q_1,\dots,q_d$, so
that $L=M\Z^d$ and $[\Z^d:L]=|\det M|$.  Since
$\conv P=M\bigl(\prod_{i=1}^d[-N_i,N_i]\bigr)$,
\begin{equation}
 \vol_d(\conv P)
 =2^d|\det M|\prod_{i=1}^dN_i
 \asymp_d[\Z^d:L]|P|,
 \label{eq:index-volume}
\end{equation}
the comparison using $N_i\ge1$ and therefore uniform in the widths.
Lemma~\ref{lem:gap-contained} and convexity give
$\conv P\subseteq C_{d,\beta}B_{\R}$, and
$\vol_d(B_{\R})\asymp_d|B|$, so equation~\eqref{eq:index-volume} implies
$[\Z^d:L]|P|\ll_{d,\beta}|B|$.  Using $|P|\ge\gamma|B|$ now gives
$[\Z^d:L]\ll_{d,\beta}\gamma^{-1}$, as claimed.
\end{proof}

\begin{lemma}[Full-dimensional projective support]
\label{lem:projective-span}
Fix $d$ and $\beta>1$, and fix a structural package arising from
Theorem~\ref{thm:cfp-structure} in ambient dimension $d$ with polynomial-box
exponent at most $\beta$.  Let $V\subseteq B\subseteq\Z^d$ be
$(\lambda,X)$-shell non-dividing and unshiftedly
$(\delta,\gamma)$-irreducible relative to this package, put $n=|V|$, and
assume
\[
 |B|\le n^\beta,
 \qquad
 [s(\delta n)]^{-1}=o(\gamma).
\]
Then, for all sufficiently large $n$,
\[
 \Span_{\R}V=\R^d,
 \qquad
 \aff\pi(V)=H_\lambda.
\]
Moreover
\[
 \pi(V)\subseteq
 \Omega:=\{y\in H_\lambda:Xy\in B_{\R}\},
\]
and $\Omega$ has nonempty relative interior in $H_\lambda$.
\end{lemma}

\begin{proof}
Apply Lemma~\ref{lem:index-bound} with $C=V$.  Since
$s(V)\ge s(\delta n)$ for all sufficiently large $n$, the hypothesis implies
$s(V)\gamma\to\infty$.  The structural step lattice $L=\langle P(V)\rangle$
therefore has rank $d$.

Let $V_0\subseteq V$ and $q$ be the structural core and translation from
Theorem~\ref{thm:cfp-structure}.  For sufficiently large $n$, every coordinate
radius of $cs(V)P(V)$ is at least one.  Since
\[
 q+cs(V)P(V)\subseteq\SigmaSet(V_0)\subseteq\Span_{\R}V,
\]
taking differences of points in the left-hand translate shows that every step
of $P(V)$ belongs to $\Span_{\R}V$.  These steps have rank $d$, so
$\Span_{\R}V=\R^d$.

If $\pi(V)$ were contained in a proper affine subspace of $H_\lambda$, then
its linear span would have dimension at most $d-1$.  Since
$v=\lambda(v)\pi(v)$ for every $v\in V$, this would put $V$ in a proper
linear subspace of $\R^d$, a contradiction.  Hence
$\aff\pi(V)=H_\lambda$.

Finally, for $v\in V$ the shell condition gives $X/\lambda(v)\in[1/2,1]$.
As $B_{\R}$ is convex, symmetric and contains the origin,
\[
 X\pi(v)=\frac{X}{\lambda(v)}v\in B_{\R}.
\]
Thus $\pi(V)\subseteq\Omega$.  Since $\Omega$ is convex and contains an
affinely spanning subset of $H_\lambda$, it has nonempty relative interior.
\end{proof}

\section{Projective convexity from a tube of subset sums}
\label{sec:projective-convexity}

We now prove the central projective statement.

\begin{definition}[Convex position for an indexed multiset]
Let $E$ be a finite-dimensional Euclidean affine space.  An indexed family
$p_1,\dots,p_n\in E$ is in $\mu$-convex position if, for every index $j$,
there is a closed half-space of $E$ containing $p_j$ and containing at most
$\mu n$ indices, counted with multiplicity.  In a zero-dimensional affine
space the only nonempty half-space is the whole space.
\end{definition}

\begin{lemma}[Projective duality]
\label{lem:projective-duality}
Let $V$ be a finite set lying in a $(\lambda,X)$-shell, and suppose that the indexed projective family $\pi(V)$ is not in $\mu$-convex position.  If $\mu n\ge4$, then there are $a\in V$ and coefficients $\beta_v$ for $v\in V\setminus\{a\}$ such that
\[
 a=\sum_{v\ne a}\beta_vv,
 \qquad
 0\le\beta_v\le\frac{4}{\mu n},
 \qquad
 \frac13\le\sum_{v\ne a}\beta_v\le3.
\]
\end{lemma}

\begin{proof}
Choose an index $a$ for which the convex-position condition fails.  Consider the compact convex set
\[
 \mathcal K=\left\{
 \sum_{v\in V}c_v(\pi(v)-\pi(a)):
 0\le c_v\le(\mu n)^{-1},\ \sum_vc_v=1
 \right\}.
\]
We claim that $0\in\mathcal K$.  Otherwise a separating functional $u$ would be negative on all of $\mathcal K$.  But failure of convex position says that the half-space
\[
 \{x:u(x-\pi(a))\ge0\}
\]
contains more than $\mu n$ indices.  Assigning total mass one, with each coefficient at most $(\mu n)^{-1}$, to those indices gives an element of $\mathcal K$ on which $u$ is nonnegative, a contradiction.

Thus
\[
 \pi(a)=\sum_vc_v\pi(v).
\]
Multiplying by $\lambda(a)$ gives
\[
 a=\sum_v c_v\frac{\lambda(a)}{\lambda(v)}v.
\]
The shell condition puts every ratio $\lambda(a)/\lambda(v)$ in $[1/2,2]$.  Move the possible $v=a$ term to the left and divide by $1-c_a$.  Since $c_a\le(\mu n)^{-1}\le1/4$, the stated bounds follow.
\end{proof}

\begin{lemma}[Robust weighted zonotope]
\label{lem:robust-zonotope}
Fix $d$ and $\beta>1$.  Let $V\subseteq B\subseteq\Z^d$, and fix a
structural package arising from Theorem~\ref{thm:cfp-structure} in ambient
dimension $d$ with polynomial-box exponent at most $\beta$, applicable to
every subset of $V$ of size at least $\delta|V|$.  Suppose that $(V,B)$ is
unshiftedly $(\delta,\gamma)$-irreducible relative to this package, put
$n=|V|$, and assume
\[
 |B|\le n^\beta,
 \qquad
 \delta\ge n^{-1/2},
 \qquad
 \delta=o(\mu),
 \qquad
 [s(\delta n)]^{-1}=o(\gamma).
\]
Let $a\in V$ and let weights $(\theta_v)_{v\in V\setminus\{a\}}$
satisfy
\[
 0\le\theta_v\le\frac12,
 \qquad
 c_1\mu n\le\sum_{v\ne a}\theta_v\le c_2\mu n
\]
for fixed positive constants $c_1,c_2$.  Let
$R\subseteq V\setminus\{a\}$ have size $o(\mu n)$, put
$C=(V\setminus\{a\})\setminus R$, and set
\[
 z_C=\sum_{v\in C}\theta_vv.
\]
Then, for a constant $c_{d,\beta,c_1,c_2}>0$,
\[
 z_C+c_{d,\beta,c_1,c_2}\mu\gamma nB_{\R}\subseteq Z(C).
\]
\end{lemma}

\begin{proof}
Let
\[
 Z^0=\sum_{v\in C}[-\theta_v,\theta_v]v.
\]
It suffices to prove $\xi nB_{\R}\subseteq Z^0$ with
$\xi=c_{d,\beta,c_1,c_2}\mu\gamma$.

Suppose $y\in\xi nB_{\R}\setminus Z^0$.  By separation, after changing the sign of the separator if necessary, there is a linear functional $u$ such that
\[
 u(y)>\sum_{v\in C}\theta_v|u(v)|.
\]
Put $z=u(y)>0$ and
\[
 H=\left\{x:\ |u(x)|<\frac{c'_d\gamma z}{\xi n}\right\}.
\]
Since $y\in\xi nB_{\R}$, this is a central slab of relative width at most $c'_d\gamma$.  Lemma~\ref{lem:slab-exclusion} therefore gives $|C\cap H|<\delta n$.

The total coefficient mass is $\asymp\mu n$.  The removal of $R$ costs $o(\mu n)$, and the mass on $H$ is at most $\delta n/2=o(\mu n)$.  Hence
\[
 \sum_{v\in C\setminus H}\theta_v\gg\mu n.
\]
For $v\notin H$,
\[
 |u(v)|\ge\frac{c'_d\gamma z}{\xi n}.
\]
Thus
\[
 \sum_{v\in C}\theta_v|u(v)|
 \gg \mu n\frac{\gamma z}{\xi n}>z
\]
when $c_{d,\beta,c_1,c_2}$ is sufficiently small, contradicting the separating
inequality.
\end{proof}

\begin{theorem}[Projective convexity]
\label{thm:projective-convexity}
Fix $d$ and $\beta>1$.  Let $V\subseteq B\subseteq\Z^d$ be an
$(\lambda,X)$-shell non-dividing set of size $n$, where $B$ is a
symmetric coordinate box and $|B|\le n^\beta$, and fix a structural
package arising from Theorem~\ref{thm:cfp-structure} in ambient dimension
$d$ with polynomial-box exponent at most $\beta$, applicable to every subset
of $V$ of size at least $\delta|V|$.  Suppose that $(V,B)$ is unshiftedly
$(\delta,\gamma)$-irreducible relative to this package.  Assume
$0<\delta,\gamma,\mu<1$ and the parameter hierarchy
\[
 \delta\ge n^{-1/2},
 \qquad
 \delta=o(\mu),
 \qquad
 \frac{n}{\log n}=o(\mu\gamma n),
 \qquad
 \gamma^{-1}=o(\mu\gamma n),
 \qquad
 [s(\delta n)]^{-1}=o(\gamma).
\]
Then the indexed family $\pi(V)$ is in $\mu$-convex position for all sufficiently large $n$.
\end{theorem}

\begin{proof}
Suppose otherwise, and use Lemma~\ref{lem:projective-duality} to obtain
$a$ and coefficients $\beta_v$.  Choose a sufficiently small absolute
constant $c_0>0$ and set
\[
 T=\lfloor c_0\mu n\rfloor,
 \qquad
 \theta_v=T\beta_v.
\]
Then
\[
 Ta=\sum_{v\ne a}\theta_vv,
 \qquad
 0\le\theta_v\le\frac12,
 \qquad
 \sum_{v\ne a}\theta_v\asymp\mu n.
\]
Put $C=V\setminus\{a\}$ and let
$\widehat C,C_0,P,q$ be the objects assigned to $C$ by the fixed structural
package.  Put
\[
 L=\langle P\rangle,
 \qquad
 \overline C=\widehat C\setminus C_0,
\]
and write $s=s(C)$.  For sufficiently large $n$, one has $|C|\ge\delta n$.  Irreducibility therefore gives $d(C)=d$ and $|P|\ge\gamma|B|$.  Moreover, the hypothesis $[s(\delta n)]^{-1}=o(\gamma)$ implies
$s(C)\gamma\to\infty$.  Since $|C|=n-1$, the bound $|B|\le n^\beta$
implies $|B|\le|C|^{2\beta}$ for sufficiently large $n$.
Lemma~\ref{lem:index-bound}, with polynomial-box exponent $2\beta$,
consequently applies before any rounding is performed and gives
\[
 \operatorname{rank}L=d,
 \qquad
 [\Z^d:L]\ll_{d,\beta}\gamma^{-1}.
\]

The discarded set $R=C\setminus\overline C$ has size
$O(n/\log n)=o(\mu n)$, so Lemma~\ref{lem:robust-zonotope} gives
\[
 \overline z+\xi_0 nB_{\R}\subseteq Z(\overline C),
 \qquad
 \overline z=\sum_{v\in\overline C}\theta_vv,
 \qquad
 \xi_0\asymp_{d,\beta}\mu\gamma,
\]
while $Ta-\overline z\in O(n/\log n)B_{\R}$.
Lemma~\ref{lem:gap-contained} gives $P\subseteq C_{d,\beta}B_{\R}$, and
$q\in sP$; hence
\[
 (Ta-\overline z)-q
 \in O\!\left(\frac{n}{\log n}+s\right)B_{\R}
 =o(\xi_0 n)B_{\R},
\]
so the choice $\xi=\xi_0/2$ gives, for all sufficiently large $n$,
\[
 (Ta-\overline z)-q+\xi nB_{\R}\subseteq\xi_0nB_{\R}.
\]
Consequently, if $y\in L\cap(Ta+\xi nB_{\R})$, then
$y-q-\overline z\in\xi_0nB_{\R}$, so $y-q\in Z(\overline C)$ and
$y\in Z(\overline C)+q+\frac{cs}{2}P$;
Lemma~\ref{lem:one-sided-filling} therefore gives the tube inclusion
\[
 \boxed{
 L\cap(Ta+\xi nB_{\R})\subseteq\SigmaSet(C).
 }
\tag{$\dagger$}
\]

Let $r$ be the order of $a+L$ in $\Z^d/L$.  Then
$r\ll_{d,\beta}\gamma^{-1}$ and $ra\in L$.  Since
$r=o(\xi n)$ and $r=o(T)$, the nearest multiple of $r$ to $T$ is positive;
call it $m_0$.  It satisfies
\[
 |m_0-T|\le\frac r2<\xi n.
\]
As $a\in B$,
\[
 m_0a\in L\cap(Ta+\xi nB_{\R}).
\]
The tube inclusion gives $m_0a\in\SigmaSet(V\setminus\{a\})$.  Applying $\lambda$ yields
\[
 m_0\lambda(a)=\sum_{v\in S}\lambda(v)
\]
for a nonempty $S\subseteq V\setminus\{a\}$, contradicting shell non-divisibility.
\end{proof}

\section{Weighted projective density increment}
\label{sec:density-increment}

We require an indexed version of the convex-density lemma.  Rather than retranscribing the geometric proof, we derive it from the published set version by a perturbation and compactness argument.

\begin{lemma}[Weighted convex-density lemma]
\label{lem:weighted-density}
Fix $r\ge1$ and $\eps>0$.  There are $\tau\in(0,1)$ and $\mu_0>0$
such that the following holds.  Let $p_1,\dots,p_n$ be an indexed family in
a convex body $\Omega\subseteq\R^r$ with nonempty interior, in
$\mu$-convex position, where $0<\mu<\mu_0$, and assume that $n$ is
sufficiently large in $r,\eps$ and $\mu$.  Then there are
\[
 \eta\in[\mu,\mu^\tau]
\]
and a compact convex set $\Omega'\subseteq\Omega$ such that
\[
 \vol_r(\Omega')\le\eta\vol_r(\Omega)
\]
and
\[
 \#\{i:p_i\in\Omega'\}
 \ge\eta^{\kappa_r+\eps}n,
 \qquad
 \kappa_r=\frac{r-1}{r+1}.
\]
For $r=1$ this means $\kappa_1=0$.
\end{lemma}

\begin{proof}
Take $\tau$ and $\mu_0$ from \cite[Lemma~1]{PhamZakharov}.  For each $j$,
choose a witnessing closed half-space $H_j$ containing $p_j$ and at most
$\mu n$ indices.  For each $m\ge1$, perturb the points successively to
distinct points $q_i^{(m)}\in H_i\cap B(p_i,1/m)$ chosen so close to
$p_i$ that $q_i^{(m)}\in H_j$ only when $p_i\in H_j$; this is possible
because only finitely many of the half-spaces $H_j$ omit $p_i$, each at
positive distance from it, and $H_i\cap B(p_i,1/m)$ is infinite since
$r\ge1$.  Then each $H_j$ contains $q_j^{(m)}$ and at most as many
perturbed points as original indices, so the set
$Q_m=\{q_i^{(m)}:1\le i\le n\}$ is in $\mu$-convex position and lies in
$\Omega_m=\Omega+m^{-1}\overline{B_2^r}$.

The set version of the Pham--Zakharov lemma gives
$\eta_m\in[\mu,\mu^\tau]$ and a convex set $K_m\subseteq\Omega_m$, which
we may take compact by passing to its closure, with
\[
 \vol_r(K_m)\le\eta_m\vol_r(\Omega_m),
 \qquad
 |Q_m\cap K_m|\ge\eta_m^{\kappa_r+\eps}n.
\]
The bodies $K_m$ lie in the fixed bounded set
$\Omega+\overline{B_2^r}$.  Pass to a subsequence on which
$\eta_m\to\eta$, the finite index set $\{i:q_i^{(m)}\in K_m\}$ is
constant, and $K_m$ converges in the Hausdorff metric to a compact convex
set $K$; Blaschke selection and the continuity of convex-body volume
under Hausdorff convergence may be found in \cite{Schneider}.  Since
$\Omega_m$ decreases to $\Omega$, one has $K\subseteq\Omega$, and every
stabilized index yields $p_i\in K$.  Consequently
\[
 \vol_r(K)\le\eta\vol_r(\Omega),
 \qquad
 \#\{i:p_i\in K\}\ge\eta^{\kappa_r+\eps}n,
\]
and $\Omega'=K$ has the required properties.
\end{proof}

We now lift a projective cap back to the lattice.

\begin{lemma}[Projective cap lifting]
\label{lem:cap-lifting}
Let $B_{\R}\subseteq\R^d$ be a symmetric coordinate box, let
$B=B_{\R}\cap\Z^d$, and let $\lambda:\R^d\to\R$ be a nonzero integer
linear functional.  Equip
\[
 H_\lambda=\{y\in\R^d:\lambda(y)=1\}
\]
with its induced Euclidean $(d-1)$-dimensional measure $m_\lambda$, and put
\[
 \Omega=\{y\in H_\lambda:Xy\in B_{\R}\}.
\]
Let $\Omega'\subseteq\Omega$ be compact and convex with
\[
 m_\lambda(\Omega')\le\eta m_\lambda(\Omega).
\]
Finally, let
\[
 W=\{v\in B:X\le\lambda(v)\le2X,\ \pi(v)\in\Omega'\}.
\]
Then there are a symmetric coordinate box $B'\subseteq\Z^r$, an
integer-linear functional $\lambda':\Z^r\to\Z$, and an injective map
$\phi:W\to B'$ such that
\[
 \lambda'(\phi(v))=\lambda(v)\qquad(v\in W),
\]
and one of the following holds:
\begin{enumerate}[label=(\roman*)]
 \item $r<d$ and $|B'|\ll_d|B|$;
 \item $r=d$ and $|B'|\ll_d\eta|B|$.
\end{enumerate}
In particular, shell non-divisibility and the shell parameter are preserved
exactly under the identification.
\end{lemma}

\begin{proof}
Write $\lambda(x)=\ell\cdot x$ for a nonzero vector $\ell\in\Z^d$.  For
a measurable set $E\subseteq H_\lambda$ and $T>0$, put
$\operatorname{Cone}_T(E)=\{ty:y\in E,\ 0\le t\le T\}$.  Integrating the
Jacobian $t^{d-1}/\|\ell\|_2$ on
$H_\lambda=\ell/\|\ell\|_2^2+\ker\lambda$ gives
\begin{equation}
 \vol_d(\operatorname{Cone}_T(E))
 =\frac{T^d}{d\|\ell\|_2}\,m_\lambda(E),
 \label{eq:perspective-volume}
\end{equation}
so the factor depending on $\lambda$ cancels from volume ratios.

Define $K=B_{\R}\cap\operatorname{Cone}_{2X}(\Omega')$; it is convex,
contains $0$, and contains $W$.  Since
$\operatorname{Cone}_{X}(\Omega)\subseteq B_{\R}$ (as $Xy\in B_{\R}$ for
$y\in\Omega$ and $B_{\R}$ is convex and contains the origin),
formula~\eqref{eq:perspective-volume} gives
\begin{align}
 \vol_d(K)
 &\le \vol_d(\operatorname{Cone}_{2X}(\Omega'))\notag\\
 &\le 2^d\eta\,
       \vol_d(\operatorname{Cone}_{X}(\Omega))
 \le2^d\eta\vol_d(B_{\R}).
 \label{eq:cap-cone-volume}
\end{align}

Put $S=(K-K)\cap\Z^d$, $H=\Span_{\R}S$ and $r=\dim H$.  Since $0\in K$
and $W\subseteq K\cap\Z^d$, one has $W\subseteq S$; also
$K-K\subseteq2B_{\R}$.  Applying Lemma~\ref{lem:convex-lattice-john} in
the lattice $H\cap\Z^d$ to the symmetric convex body $(K-K)\cap H$ gives
a proper symmetric homogeneous rank-$r$ GAP $Q\subseteq H\cap\Z^d$ such
that
\[
 S\subseteq Q,
 \qquad
 |Q|\ll_d|S|.
 \tag{*}
\]
If $r<d$, then, writing $B_{\R}=\prod_i[-M_i,M_i]$,
\[
 |S|
 \le |2B_{\R}\cap\Z^d|
 =\prod_{i=1}^d(4M_i+1)
 \le2^d|B|,
\]
so $(*)$ gives $|Q|\ll_d|B|$.  If $r=d$, then $K-K$, and hence $K$, is
full-dimensional, and the full-rank part of
Lemma~\ref{lem:convex-lattice-john}, the Rogers--Shephard
inequality~\cite{Schneider} and \eqref{eq:cap-cone-volume} give
\[
 |Q|\ll_d|S|
 \ll_d\vol_d(K-K)
 \le \binom{2d}{d}\vol_d(K)
 \ll_d\eta\vol_d(B_{\R})
 \ll_d\eta|B|.
\]

Finally write $Q=\bigl\{\sum_{i=1}^r n_iq_i:|n_i|\le N_i\bigr\}$ and put
\[
 B'=\prod_{i=1}^r[-N_i,N_i]\cap\Z^r,
 \qquad
 \phi=\phi_Q|_W,
 \qquad
 \lambda'(n_1,\dots,n_r)=\sum_{i=1}^r n_i\lambda(q_i),
\]
where $\phi_Q$ is the coefficient identification of the proper GAP $Q$;
then $|B'|=|Q|$ and $\phi$ is injective.
Lemma~\ref{lem:functional-transform} gives
$\lambda'(\phi(v))=\lambda(v)$ for $v\in W$.  The preceding bounds for
$|Q|$ give alternatives (i) and (ii), and the identity of functional
values proves the final assertion.
\end{proof}

\section{The density-increment iteration}
\label{sec:iteration}

We use the exponent sequence defined in
Section~\ref{sec:projective-framework}.

\begin{lemma}[Uniform up-move numerology]
\label{lem:numerology}
Fix $D\ge2$ and $\zeta_0>0$.  There are $\eps_0,c_*>0$ such that, whenever
\[
 1\le d<e\le D,
 \qquad
 \zeta_0\le\zeta\le1-\sigma_d,
 \qquad
 0<\eps<\eps_0,
\]
one has
\[
 (\sigma_e+\zeta)
 \left(\frac1{\sigma_d+\zeta}-(1-\eps)(e-d)\right)
 \le1-c_*
\]
whenever the parenthetical expression is nonnegative.
\end{lemma}

\begin{proof}
For $k=e-d$, set
\[
 f_{d,k}(\zeta)
 =(\sigma_{d+k}+\zeta)
  \left(\frac1{\sigma_d+\zeta}-k\right).
\]
At $\zeta=0$, direct substitution gives $f_{d,k}(0)\le1$; equality occurs
only for
\[
 (d,e)\in\{(1,2),(1,3),(1,4),(2,3),(2,4),(3,4)\}.
\]
Moreover
\[
 f_{d,k}'(\zeta)
 =\frac{\sigma_d-\sigma_{d+k}}{(\sigma_d+\zeta)^2}-k<0.
\]
Hence the required gap is strictly positive at $\eps=0$ for every pair and
every $\zeta\in[\zeta_0,1-\sigma_d]$ for which the parenthetical factor is
nonnegative.  There are finitely many pairs $1\le d<e\le D$; compactness and
continuity give a uniform gap $c_*>0$, which persists after choosing
$\eps_0>0$ sufficiently small.
\end{proof}

\begin{theorem}[Projective theorem]
\label{thm:projective-main}
For every fixed $d$ and $\beta>1$, every $(\lambda,X)$-shell
non-dividing set $V\subseteq B\subseteq\Z^d$, where $B$ is a symmetric
coordinate box and $|B|\le|V|^\beta$, satisfies
\[
 |V|\le |B|^{\sigma_d+o(1)}.
\]
Here the $o(1)$ is uniform over the shell functional and shell parameter once
$d$ and $\beta$ are fixed.
\end{theorem}

\begin{proof}
Assume for a contradiction that for some fixed $\zeta_0>0$ there are
arbitrarily large examples
\[
 A_0\subseteq B_0\subseteq\Z^{d_0},
 \qquad
 |A_0|>|B_0|^{\sigma_{d_0}+\zeta_0}.
\]
Write $n_0=|A_0|$.  Since $|A_0|\le|B_0|$, the assumed examples are
already impossible when $\sigma_{d_0}+\zeta_0\ge1$; hence we may and do
assume $\sigma_{d_0}+\zeta_0<1$.  All sets constructed below remain shell
non-dividing, because the only operations are passage to subsets and
coefficient identifications of proper homogeneous GAPs, with the shell
functional transported by Lemma~\ref{lem:functional-transform}.

\medskip\noindent\emph{Choice of parameters.}\quad At any
noncontradictory stage $A\subseteq B\subseteq\Z^d$ with excess at least
$\zeta_0$, the inequalities $|A|\le|B|$ and
$|A|>|B|^{\sigma_d+\zeta_0}$ give
\[
 |B|<|A|^{\beta_*},
 \qquad
 \beta_*=(\sigma_1+\zeta_0)^{-1}>1.
\]
Thus one fixed polynomial-box exponent works at every stage.  Since
$\sigma_r\to1$, the integer
\[
 D_0=\max\{r\ge1:\sigma_r\le1-\zeta_0\}
\]
is finite.  At the beginning of every noncontradictory stage, the exact
excess is at least $\zeta_0$ and $|A|\le|B|$, so its ambient dimension is
at most $D_0$.

For each $1\le\ell\le D_0$, let $G_\ell$ be a bound for the output
dimension in Lemma~\ref{lem:unshifted-preprocessing} with ambient dimension
$\ell$ and polynomial-box exponent $\beta_*$.  Choose
\[
 D>\max\{D_0,G_1,\dots,G_{D_0}\}.
\]
Thus every one-step preprocessing output has dimension at most $D$, while
every stage that survives to the next iteration has dimension at most $D_0$.
Apply Lemma~\ref{lem:numerology} with $D$ and $\zeta_0$,
obtaining a uniform up-move gap $c_*>0$.

Put
\[
 \Delta_D=\min_{2\le r\le D}(\sigma_r-\sigma_{r-1})>0.
\]
Choose $\eps>0$ so small that
\[
 \eps<\min\left\{
  \eps_0,\frac{\zeta_0}{4},\frac{c_*}{4},
  \frac{\Delta_D}{8},\frac13
 \right\}
\]
and so that $\kappa_{r-1}+\eps<\sigma_r+\zeta_0$ for every
$2\le r\le D$.  Choose $K_0$ sufficiently large that
\[
 \frac1{K_0}<\sigma_1+\zeta_0,
 \qquad
 K_0(\sigma_1+\zeta_0)>2.
\]
For $1\le\ell\le D_0$, let
$K_\ell=K(\ell,\beta_*,\eps,K_0)$ be an admissible exponent in
Lemma~\ref{lem:unshifted-preprocessing}, and put
\[
 K=\max_{1\le\ell\le D_0}K_\ell.
\]
Using this larger exponent is valid in every ambient dimension at most
$D_0$.  By the final assertion of
Lemma~\ref{lem:unshifted-preprocessing}, there is a fixed number
$\beta_{\rm pkg}>1$, depending only on the constants chosen so far, which
bounds the polynomial-box exponent of every structural package output by
preprocessing in dimensions at most $D_0$.  Put
\[
 \beta^\dagger
 =\max\left\{\beta_{\rm pkg},
       \frac{2\beta_*}{1-\eps}\right\}.
\]
Finally choose $0<\kappa<1$ sufficiently small in the fixed parameters.

We now choose a slowly varying scale.  For each $1\le r\le D-1$,
let $N_r(t)$ be a threshold for Lemma~\ref{lem:weighted-density} in
dimension $r$ at parameter $\mu=t$.  A diagonal choice gives a nondecreasing
integer-valued function $\omega(n)\to\infty$ such that
\[
 \omega(n)^{2K+2}=o(\log n),
 \qquad
 N_r(\omega(n)^{-\kappa})\le n^{1-\eps}
 \quad(1\le r\le D-1).
\]
For completeness, choose integers $L_m\uparrow\infty$ so that
$\log L_m\ge m^{4K+4}$ and
$L_m^{1-\eps}\ge\max_rN_r(m^{-\kappa})$, and put
$\omega(n)=\max\{m:L_m\le n\}$ (with any fixed initial value before
$L_2$).  No quantitative dependence of $N_r(t)$ on $t$ is required.

At a stage with $n=|A|$, use
\[
 \delta=\omega(n)^{-1},
 \qquad
 \gamma=\delta^K,
 \qquad
 \mu=\delta^\kappa.
\]
Uniformly as $n\to\infty$, these parameters satisfy
\begin{align*}
 \delta&\ge n^{-1/2},&
 \gamma&\ge n^{-1/3},&
 \delta&=o(\mu),\\
 \frac1{\log n}&=o(\mu\gamma),&
 \gamma^{-1}&=o(\mu\gamma n),&
 [s(\delta n)]^{-1}&=o(\gamma).
\end{align*}
Each condition follows from $\omega(n)\to\infty$ together with the growth
condition $\omega(n)^{2K+2}=o(\log n)$, which also forces
$\omega(n)=n^{o(1)}$; we omit the routine verification, which is carried
out in full in the formalization.  The parameters therefore meet all
hypotheses of the preprocessing theorem, and the last condition in the
definition of $\omega$ will supply the largeness hypothesis in the
weighted convex-density lemma after preprocessing.

\medskip\noindent\emph{The one-step increment.}\quad Suppose
$A\subseteq B\subseteq\Z^d$ is a current stage, put $n=|A|$, and define
its exact excess by
\[
 p=\frac{\log n}{\log|B|}=\sigma_d+\zeta.
\]
The case $|B|=1$ is already impossible for large $n$.  Since $n\le|B|$,
one has
\[
 \zeta_0\le\zeta\le1-\sigma_d.
\]
In particular $\sigma_d\le1-\zeta_0$, so $d\le D_0$.  Apply
Lemma~\ref{lem:unshifted-preprocessing} with the ambient-dimension parameter
$d$, and denote its output by
$\widetilde A\subseteq\widetilde B\subseteq\Z^e$.  Put
$\widetilde n=|\widetilde A|$.  Since
$\widetilde n>n^{1-\eps}$ and $\log\widetilde n\asymp\log n$, the same
parameters satisfy, with $\widetilde n$ in place of $n$,
\[
 \delta\ge\widetilde n^{-1/2},
 \qquad
 \frac1{\log\widetilde n}=o(\mu\gamma),
 \qquad
 \gamma^{-1}=o(\mu\gamma\widetilde n),
 \qquad
 [s(\delta\widetilde n)]^{-1}=o(\gamma).
\]
In every preprocessing alternative one has
$|\widetilde B|\ll n^{\beta_*}$; since
$\widetilde n>n^{1-\eps}$, it follows, after absorbing the fixed implicit
constant, that
\[
 |\widetilde B|\le \widetilde n^{\beta^\dagger}
\]
for all sufficiently large $n$.  By the definition of $\beta^\dagger$, the retained structural package
arises with polynomial-box exponent at most $\beta^\dagger$.  Thus the same
package is used in every subsequent
application of projective convexity, slab exclusion, filling and the index
lemma.  Finally,
\[
 \widetilde n>n^{1-\eps}
 \ge N_r(\mu)\qquad(1\le r\le D-1),
\]
by the defining property of $\omega$.  Thus every subsequent use of
projective convexity and weighted convex density is legitimate for the
post-preprocessing set.

If $e>d$, put
\[
 a=\frac1p-(1-\eps)(e-d).
\]
The box estimate gives $|\widetilde B|\ll n^a$.  If $a\le0$, this is
impossible for large $n$.  Otherwise Lemma~\ref{lem:numerology}
gives
\[
 (\sigma_e+\zeta)a\le1-c_*.
\]
Since $a\le\beta_*$, after decreasing a constant
$c_{\rm up}=c_{\rm up}(D,\zeta_0)>0$ and absorbing the fixed implicit
constant, one obtains
\[
 |\widetilde A|>n^{1-\eps}
 >|\widetilde B|^{\sigma_e+\zeta+c_{\rm up}}.
\]
Thus an upward dimension change increases the excess by at least the fixed
amount $c_{\rm up}$.  Any noncontradictory output satisfies
$\sigma_e+\zeta+c_{\rm up}<1$, and hence $e\le D_0$; the stage-dimension
invariant is therefore preserved.

If $e<d$, then
\[
 |\widetilde A|>n^{1-\eps},
 \qquad
 |\widetilde B|\ll|B|.
\]
Since $p\le1$ and $\sigma_d-\sigma_e\ge\Delta_D$, the choices above give,
after absorbing the fixed constant,
\[
 |\widetilde A|
 >|\widetilde B|^{\sigma_e+\zeta+c_{\rm down}}
\]
for a fixed $c_{\rm down}>0$ depending only on $D$ and $\zeta_0$.

It remains to consider $e=d$.  Write
\[
 |\widetilde A|=\rho n,
 \qquad
 |\widetilde B|\ll\rho^{K_0}|B|.
\]
The pair $(\widetilde A,\widetilde B)$ is irreducible.  Let
$\widetilde\lambda$ and $X$ be the transported shell functional and shell
parameter.  If $d=1$, Theorem~\ref{thm:projective-convexity} says that its
projective indexed family is in $\mu$-convex position.  It lies in the
zero-dimensional affine space $H_{\widetilde\lambda}$, whose only nonempty
half-space is the whole space; that half-space contains every index.  Since
$\mu<1$, this is impossible.

Assume $d\ge2$.  By Lemma~\ref{lem:projective-span}, the
family $\pi(\widetilde A)$ affinely spans $H_{\widetilde\lambda}$ and lies
in the convex body
\[
 \Omega=\{y\in H_{\widetilde\lambda}:Xy\in\widetilde B_{\R}\},
\]
which has nonempty relative interior.  Theorem~\ref{thm:projective-convexity}
and Lemma~\ref{lem:weighted-density}, applied in the Euclidean affine space
$H_{\widetilde\lambda}\cong\R^{d-1}$, followed by
Lemma~\ref{lem:cap-lifting}, give a shell non-dividing set $A'$ with
\[
 |A'|\ge\rho\eta^{a_d}n,
 \qquad
 a_d=\kappa_{d-1}+\eps,
 \qquad
 \eta\in[\mu,\mu^{\widehat\tau_d}],
\]
where
\[
 \widehat\tau_d:=\tau_{d-1}>0
\]
depends only on $d$ and $\eps$.  If cap lifting lowers the
dimension to some $f<d$, then its box $B'$ satisfies
$|B'|\ll|\widetilde B|\ll|B|$.  Moreover
$\eta^{a_d}=n^{-o(1)}$, and hence for sufficiently large $n$,
\[
 |A'|\ge n^{1-2\eps}.
\]
The same calculation as in the preceding down case gives
\[
 |A'|>|B'|^{\sigma_f+\zeta+c_{\rm down}}.
\]
Thus every dimension-lowering cap also produces a fixed excess increment.

Suppose instead that cap lifting preserves dimension $d$.  Then
\[
 |B'|\ll\eta\rho^{K_0}|B|.
\]
At the old exponent $p=\sigma_d+\zeta$ we have
\[
 \frac{|A'|}{|B'|^p}
 \gg
 \eta^{a_d-p}\rho^{1-K_0p}.
\]
The choices of $\eps$ and $K_0$ give, uniformly for $2\le d\le D$,
\[
 p-a_d\ge\nu_0>0,
 \qquad
 K_0p-1\ge1.
\]
Let
\[
 \tau_* =\min_{2\le r\le D}\widehat\tau_r>0.
\]
Since $\eta\le\mu^{\widehat\tau_d}$ and $\rho\le1$, the right-hand side is at least
\[
 C_D^{-1}\mu^{-\tau_*\nu_0},
\]
which tends to infinity.  Therefore, for all sufficiently large $n$,
\[
 |A'|>|B'|^p.
\]
The exact excess strictly increases in a same-dimensional step.  In
addition,
\[
 \frac{|B'|}{|B|}
 \ll_D\eta\rho^{K_0}
 \le C_D\mu^{\tau_*}=o(1).
\]
This completes the one-step analysis.

\medskip\noindent\emph{The iteration terminates.}\quad We first prove
that the iteration cannot lose too many points.  An upward
step increases the excess by at least $c_{\rm up}$, while the exact excess
is always at most $1-\sigma_d<1$.  Hence there are only
$O_{D,\zeta_0}(1)$ upward steps.  Every downward step decreases the
integer dimension, and between two upward steps there can be at most $D$
downward steps.  Since all dimensions are at most $D$, the total number
of dimension-changing steps is at most some fixed number
$M=M(D,\zeta_0)$.

For $2\le r\le D$ put
\[
 c_r=\max\left\{\kappa_{r-1}+\eps,\frac1{K_0}\right\}.
\]
By construction
\[
 c_r<\sigma_r+\zeta_0.
\]
Consider any consecutive block of same-dimensional steps in dimension
$r$.  At a step in this block, let $\rho_j$ be the preprocessing
cardinality ratio and let $\eta_j$ be the cap-volume ratio.  Put
\[
 x_j=\log(1/\eta_j),
 \qquad
 y_j=\log(1/\rho_j).
\]
The point and box estimates give
\begin{align*}
 \log\frac{|A_j|}{|A_{j+1}|}
 &\le (\kappa_{r-1}+\eps)x_j+y_j,\\
 \log\frac{|B_j|}{|B_{j+1}|}
 &\ge x_j+K_0y_j-O_D(1).
\end{align*}
The first right-hand side is at most
$c_r(x_j+K_0y_j)$.  Also
\[
 x_j\ge\tau_*\log(1/\mu_j)
 =\tau_*\kappa\log\omega(|A_j|)\longrightarrow\infty
\]
uniformly whenever all current cardinalities are bounded below by a fixed
positive power of $n_0$, because $\omega$ is nondecreasing and tends to
infinity.  Thus the $O_D(1)$ term is
$o(x_j+K_0y_j)$ uniformly.  Summing over any initial segment of the block
gives
\[
 \log\frac{|A_{\rm start}|}{|A_{\rm end}|}
 \le(c_r+o(1))\log|B_{\rm start}|.
\]
Since the starting excess is at least $\zeta_0$,
\[
 |A_{\rm end}|
 \ge
 |A_{\rm start}|^{1-c_r/(\sigma_r+\zeta_0)-o(1)}.
\]
Therefore there is a fixed $\theta>0$, depending only on the chosen
parameters, such that every same-dimensional block, including every
partial block, retains at least the $\theta$-th power of its input
cardinality.  Every dimension-changing step retains at least
$|A_j|^{1-2\eps}$ points.

There are at most $M$ dimension-changing steps and at most $M+1$
same-dimensional blocks.  Hence there is a fixed $q>0$ such that every
iterate produced before a contradiction satisfies
\[
 |A_j|\ge n_0^q.
\]
To remove the apparent circularity in the uniform $o(1)$ above, run the
construction with the provisional stopping rule $|A_j|<n_0^{q/2}$.  On
the range before this rule triggers, all $o(1)$ estimates are uniform;
the preceding product calculation yields the stronger lower bound
$|A_j|\ge n_0^q$, so the provisional rule never triggers.

We may therefore iterate indefinitely unless a contradiction occurs.
There are only finitely many dimension changes.  After the last one, every
step has the same dimension, and the box estimate above gives, uniformly
for all remaining stages and sufficiently large $n_0$,
\[
 |B_{j+1}|\le\frac12|B_j|.
\]
A positive integer box cardinality cannot be halved indefinitely, while
the accompanying sets all have cardinality at least $n_0^q>1$.  This is
the required contradiction.
\end{proof}

\section{Proof of the main theorem}
\label{sec:main-proof}

\begin{proof}[Proof of Theorem~\ref{thm:main}]
The lower bound follows from Proposition~\ref{prop:lower-bound}.  For the
upper bound, fix $0<\varepsilon<1$ and suppose, for arbitrarily large $N$, that
there is a non-dividing set $A\subseteq[N]$ with
\[
 |A|>N^{1/5+2\varepsilon}.
\]
Partition $[N]$ into dyadic intervals.  Some interval $[X,2X)$ contains a
subset $A_0\subseteq A$ with
\[
 |A_0|\ge\frac{|A|}{1+\log_2N}>N^{1/5+\varepsilon}
\]
for all sufficiently large $N$.  Take
\[
 d=1,
 \qquad
 B=[-2X,2X]\cap\Z,
 \qquad
 \lambda(t)=t.
\]
Then $A_0$ is shell non-dividing and $|B|\le4N+1$.  Consequently, for all
sufficiently large $N$,
\[
 |A_0|>|B|^{1/5+\varepsilon/2},
 \qquad
 |B|<|A_0|^{\beta},
 \qquad
 \beta=(1/5+\varepsilon/2)^{-1}>1.
\]
Theorem~\ref{thm:projective-main}, applied with this fixed $\beta$, gives
\[
 |A_0|\le |B|^{1/5+o(1)}.
\]
For sufficiently large $N$ the $o(1)$ term is smaller than
$\varepsilon/2$, contradicting the preceding strict inequality.  Hence
$F(N)\le N^{1/5+o(1)}$, and the theorem follows.
\end{proof}

\section{A concluding remark}
\label{sec:remarks}

For $x=\min A$, a residue fiber $A_r=\{a\in A:a\equiv r\pmod x\}$ of a
non-dividing set $A\subseteq[N]$ becomes a non-averaging set after the
affine normalization $a\mapsto(a-r)/x$.  If a fiber has size $h$, the
sharp non-averaging bound therefore gives $N\ge xh^{4-o(1)}$, so a direct
multiplicity amplification estimate of the form
\[
 h\ge |A|^{5/4-o(1)}x^{-1/4}
\]
would give another proof of Theorem~\ref{thm:main}.  Such a statement is
false for arbitrary zero-sum-free sequences, so it would have to exploit
the simultaneous forbidden-modulus property.

\begin{thebibliography}{99}

\bibitem{ErdosProblems}
T.~F.~Bloom,
\emph{Erd\H{o}s problems}, Problem~131,
\url{https://www.erdosproblems.com/131}.

\bibitem{CFPHomogeneous}
D.~Conlon, J.~Fox and H.~T.~Pham,
\emph{Homogeneous structures in subset sums and non-averaging sets},
to appear in Commun. Amer. Math. Soc.; arXiv:2311.01416.

\bibitem{Er97}
P.~Erd\H{o}s,
\emph{Some old and new problems in various branches of combinatorics},
Discrete Math. \textbf{165--166} (1997), 227--231,
\href{https://doi.org/10.1016/S0012-365X(96)00173-2}{doi:10.1016/S0012-365X(96)00173-2}.

\bibitem{ELRSS99}
P.~Erd\H{o}s, V.~Lev, G.~Rauzy, C.~S\'andor and A.~S\'ark\H{o}zy,
\emph{Greedy algorithm, arithmetic progressions, subset sums and divisibility},
Discrete Math. \textbf{200} (1999), 119--135,
\href{https://doi.org/10.1016/S0012-365X(98)00385-9}{doi:10.1016/S0012-365X(98)00385-9}.

\bibitem{PhamZakharov}
H.~T.~Pham and D.~Zakharov,
\emph{Sharp bound for the Erd\H{o}s--Straus non-averaging set problem},
Geom. Funct. Anal. \textbf{35} (2025), 1712--1738,
\href{https://doi.org/10.1007/s00039-025-00728-8}{doi:10.1007/s00039-025-00728-8}.

\bibitem{Schneider}
R.~Schneider,
\emph{Convex Bodies: The Brunn--Minkowski Theory},
2nd expanded ed., Encyclopedia of Mathematics and its Applications 151,
Cambridge University Press, Cambridge, 2014.

% For double-blind submission, replace the repository URL by an
% anonymized artifact link (e.g. anonymous.4open.science or the journal's
% ancillary-files mechanism).
\bibitem{Formalization}
Lean~4 formalization of the present paper,
\url{https://github.com/theofilxeff/erdos_131}.

\end{thebibliography}

\end{document}
